
For box QPs, the next lemma certifies growth from the gradient and the
Hessian at a point, and shows what minimality and growth imply for the
Hessian. For a point $x^*\in X$ let
$J_0=\{i\in I_C:\ell_i<x^*_i<u_i\}$ be its continuous coordinates strictly
inside their bounds and $J_\partial=\{i:x^*_i\in\{\ell_i,u_i\}\}$ its
coordinates at a bound.

\begin{lemma}[Growth certificates for box QPs]\label{lem:growthcert}
Let $F=\frac12x^\T Hx+b^\T x+c$ be a quadratic on a mixed box $X$, $x^*\in X$
and $\zeta=\nabla F(x^*)$.
\begin{enumerate}
\item[(a)] Suppose that $\zeta_i=0$ for $i\in J_0$, and that for every
$i\in I_C\setminus J_0$ the gradient points inward: $\zeta_i\ge0$ if
$x^*_i=\ell_i$ and $\zeta_i\le0$ if $x^*_i=u_i$. Let $M=\diag(\mu)$ with
$\mu_i=0$ for $i\in J_0$, $\mu_i=\abs{\zeta_i}/s_i$ for $i\in J_\partial$
with inward $\zeta_i$, and $\mu_i=-\abs{\zeta_i}$ for the remaining
coordinates, which are integer coordinates. Then
\begin{equation}\label{eq:minorant}
F(x)-F(x^*)\ge\tfrac12(x-x^*)^{\T}(H+2M)(x-x^*)\qquad(x\in X).
\end{equation}
Consequently, if $H+2M\succeq2g_0I$ for some $g_0>0$, then $x^*$ is the unique
minimizer and $F$ has point growth with constant $g_0$.
\item[(b)] Conversely, let $x^*$ be a minimizer. Then $\zeta_{J_0}=0$ and
$H_{J_0J_0}\succeq0$. If $F$ has point growth with constant $g$, then
$H_{J_0J_0}\succeq2gI$; if it has weighted growth with constant $\gamma$ at
$x^*$, then $H_{J_0J_0}\succeq2\gamma\diag(L_i)_{i\in J_0}$. In particular,
under point growth $g\le\frac12v^{\T}Hv/\norm v^2$ for every nonzero $v$
supported on $J_0$, and if $J_0\ne\emptyset$, every valid curvature bound
satisfies $L\ge2g$, so $\kappa\ge2$. If $J_0=[n]$, then $F$ is convex, and
strongly convex under point growth.
\end{enumerate}
\end{lemma}

\begin{proof}
(a) Let $x\in X$ and $d=x-x^*$. Since $F$ is quadratic,
$F(x)-F(x^*)=\zeta^{\T}d+\frac12d^{\T}Hd$ exactly, and it suffices to show
$\zeta_id_i\ge\mu_id_i^2$ for every $i$. For $i\in J_0$ both sides vanish. If
$x_i^*=\ell_i$ and $\zeta_i\ge0$, then $0\le d_i\le s_i$, so
$\zeta_id_i=\abs{\zeta_i}d_i\ge\abs{\zeta_i}d_i^2/s_i=\mu_id_i^2$; the case
$x_i^*=u_i$, $\zeta_i\le0$ is symmetric. For the remaining $i\in I_Z$, $d_i\in\Z$
and $\abs{d_i}\le d_i^2$, which gives
$\zeta_id_i\ge-\abs{\zeta_i}\abs{d_i}\ge-\abs{\zeta_i}d_i^2=\mu_id_i^2$. If
$H+2M\succeq2g_0I$, the right side of~\eqref{eq:minorant} is at least
$g_0\norm d^2$, which is positive for $d\ne0$.

(b) For $i\in J_0$ both directions $\pm e_i$ are feasible from $x^*$, so
minimality gives $\zeta_i=0$. Let $v\ne0$ be supported on $J_0$. For all
sufficiently small $\abs t$, $x^*+tv\in X$ and
$F(x^*+tv)-\OPT=\frac{t^2}2v^{\T}Hv$. This is nonnegative by minimality;
point growth bounds it from below by $gt^2\norm v^2$, and weighted growth by
$\gamma t^2\norm v_{\mathbf L}^2=\gamma t^2\sum_{i\in J_0}L_iv_i^2$; this gives
the three matrix inequalities. Taking $v=e_i$ with $i\in J_0$ gives
$H_{ii}\ge2g$ under point growth, and every valid curvature bound satisfies
$L\ge L_i\ge H_{ii}$, because $t\mapsto F(x+te_i)-L_it^2/2$ has second
derivative $H_{ii}-L_i$. If $J_0=[n]$, then $H=H_{J_0J_0}$.
\end{proof}

\begin{remark}[What growth restricts]\label{rem:nonconvex}
For a box QP, the block $H_{J_0J_0}$ at a minimizer is positive semidefinite
(Lemma~\ref{lem:growthcert}(b)), so at every minimizer, with or without
growth, every direction of negative curvature involves active bounds or
integer coordinates. Point growth with some $g>0$ is equivalent to
uniqueness of the minimizer (Lemma~\ref{lem:unique-growth}), so what
restricts the instances is the size of the constants. Point growth gives
$H_{J_0J_0}\succeq2gI$, hence $\kappa\ge2L/\lambda_{\min}(H_{J_0J_0})$ if
$J_0\ne\emptyset$, and weighted growth gives
$H_{J_0J_0}\succeq2\gamma\diag(L_i)_{i\in J_0}$. For every $x'\in X$ with
$F(x')>\OPT$, \eqref{eq:growth} gives
$\kappa\ge L\norm{x'-x^*}^2/(F(x')-\OPT)$ and \eqref{eq:wgrowth} gives
$\bar\kappa\ge\norm{x'-x^*}_{\mathbf L}^2/(F(x')-\OPT)$. Since
$\norm d_{\mathbf L}^2\le L\norm d^2$, nearly optimal points far from $x^*$
in the norm $\norm\cdot_{\mathbf L}$ make both condition numbers large.
Within these conditions an instance can have an indefinite Hessian,
$\kappa\le\frac{21}4$ and exponentially many strict local minima
(Example~\ref{ex:family}).
Part~(a) certifies growth when the first-order terms $\zeta_id_i$ at active
bounds supply the growth that the Hessian lacks. At integer coordinates it
uses integrality only to bound $\zeta_id_i$ below by $-\abs{\zeta_i}d_i^2$.
\end{remark}
